\chapter{Introduction}
\label{ch:introduction}

% --------------------------------------------------------------------------
\section{Motivation}
\label{sec:motivation}
% --------------------------------------------------------------------------

\TODO{%
  Introduce eBPF: general-purpose kernel extension framework used by Cilium,
  Falco, Tetragon, XDP. Programs run inside the kernel at near-native speed.
  The verifier is the last safety gate before a program enters the kernel.

  One-paragraph supply-chain hook: a single verifier bug allows an adversary
  to escape the sandbox at full kernel privilege. Cite CVE-2023-2163,
  CVE-2024-41003, and SoK S\&P 2025~\cite{sok-ebpf} for ``100+ verifier bugs
  reported in 2024''. Then drop the supply-chain framing --- it appears in this
  section only, not again.

  Close with why fuzzing is hard: programs must pass structural validity checks
  before the verifier even runs them; random bytecode is nearly always rejected
  before any interesting code path is reached.
}

% --------------------------------------------------------------------------
\section{Problem Statement}
\label{sec:problem}
% --------------------------------------------------------------------------

\TODO{%
  State the gap precisely: no existing system combines LLM-guided eBPF program
  generation with live KCOV kernel PC coverage as a reinforcement-learning
  feedback signal for the BPF verifier.

  Existing eBPF fuzzers (buzzer~\cite{buzzer}, BRF~\cite{brf}) do not use LLM
  generation and buzzer's KCOV integration is display-only (not feedback-driven).
  Existing LLM fuzzers (CovRL~\cite{covrl}, Fuzz4All~\cite{fuzz4all}) target
  JavaScript engines or C programs, not the BPF verifier with live kernel instrumentation.

  This work fills the gap.
}

% --------------------------------------------------------------------------
\section{Contributions}
\label{sec:contributions}
% --------------------------------------------------------------------------

This thesis makes the following contributions:

\begin{enumerate}

  \item \textbf{End-to-end pipeline.}
        An SFT + GRPO pipeline that drives a language model with live KCOV kernel
        PC coverage as an RL reward signal for eBPF verifier exploration.
        The pipeline is built on top of a modified buzzer data collector and a
        KCOV-instrumented Linux~6.8 VM.

  \item \textbf{SFT result.}
        QLoRA fine-tuning of Qwen2.5-Coder-1.5B raises the BPF verifier
        pass-rate from 1\% (zero-shot base model) to 60\%, demonstrating that a
        1.5B-parameter model can learn BPF verifier syntax from 27\,k labeled
        programs.

  \item \textbf{Failure-mode diagnosis.}
        The RL training collapse --- \texttt{reward\_std}${}={}$0 at step
        $\approx$1300 of 8370 --- is traced to a \texttt{kcov\_validator} bug
        that discarded the KCOV PC trace for rejected programs.  All rejected
        programs appeared identical to GRPO, zeroing the gradient signal.  The
        diagnosis is grounded in the GRPO gradient-starvation literature
        \cite{grpo-starvation}.

  \item \textbf{Reward redesign.}
        A depth-based, verdict-blind reward function is proposed and implemented:
        programs are scored by verifier-path depth and a discovery bonus for
        new kernel PCs, independent of accept/reject verdict.
        Run~2 (future work, Chapter~\ref{ch:conclusions}) will validate this redesign.

\end{enumerate}

% --------------------------------------------------------------------------
\section{Thesis Structure}
\label{sec:structure}
% --------------------------------------------------------------------------

\TODO{%
  Write this section last, once all chapters exist.
  One sentence per chapter:
  Chapter~2 provides background on eBPF, KCOV, coverage-guided fuzzing,
    LLM fine-tuning, and GRPO.
  Chapter~3 surveys related work in three clusters: eBPF-specific fuzzers,
    LLM-based fuzzing, and RL for code generation.
  Chapter~4 describes the methodology, including the buzzer-KCOV discovery
    and both reward designs.
  Chapter~5 details the implementation of the BPF encoder, kcov\_validator,
    training loop, and remote reward pipeline.
  Chapter~6 presents the experimental results: SFT pass-rate, RL run~1
    coverage exploration, plateau diagnosis, and reward redesign validation.
  Chapter~7 draws conclusions and outlines future work.
}
